\documentclass[14pt,aspectratio=1610]{beamer}
%\documentclass[14pt,aspectratio=1610,handout]{beamer}
\usepackage{csu-theme}
\usepackage{binary-trees}
\usepackage{booktabs}

\title{\Huge{} Binary Search Trees\\Intro to Traversal}
\subtitle{CSCI 315: \textit{Data Structure Analysis}\\Chapter 19}
%\author{Sean T. Hayes, Ph.D.}
\institute{Charleston Southern University}

%% Comment out date to use default (today)
% \date{March 4, 2025}
\subject{Software Programming}
%\keywords{}
\graphicspath{{../imgs/binary-tree}}

\setlength{\parskip}{0.5em}

% Don't show circle around leaves (for null pointers)
\forestset{
  default preamble={%
    for leaves={%
      justText}
  },
}

\begin{document}

\begin{frame}
  \titlepage{}
\end{frame}

\section*{Traversal}

\begin{frame}{Binary Search Tree Traversal}
  What and why:
  \begin{itemize}
    \item Visit every node in a BST exactly once.
    \item Common orders: in-order, pre-order, post-order (depth-first); level-order (breadth-first).
    \item Uses: produce sorted output, copy or delete a tree, compute aggregates (size, sum, height).
  \end{itemize}
\end{frame}

\begin{frame}{Traversing the Entire Tree}
  \emph{Traversing} a tree means ``\emph{visiting}'' each node exactly once.
  
  We may {visit} each the nodes in a tree to determine \\
  its size, display it, update it, delete it, sum the values, etc.
\end{frame}


\begin{frame}{Traversing the Entire Tree}
  \emph{Traversing} a tree means ``\emph{visiting}'' each node exactly once.
  
  We may {visit} each the nodes in a tree to determine \\
  its size, display it, update it, delete it, sum the values, etc.
\end{frame}
\begin{frame}{We can visit the notes in various orders.}
\begin{columns}
\onslide<2->{%
  \column{0.3\textwidth}
  In-order
  \begin{enumerate}
  \item
    Traverse the left subtree.
  \item
    \textbf{\color{CSUcyan}Visit the node.}
  \item
    Traverse the right subtree.
  \end{enumerate}%
}
\onslide<3->{%
  \column{0.3\textwidth}
  Pre-order
  \begin{enumerate}
  \item
    \textbf{\color{CSUcyan}Visit the node.}
  \item
    Traverse the left subtree.
  \item
    Traverse the right subtree.
  \end{enumerate}%
}
\onslide<4->{%
\column{0.3\textwidth}
  Post-order
  \begin{enumerate}
  \item
    Traverse the left subtree.
  \item
    Traverse the right subtree.
  \item
    \textbf{\color{CSUcyan}Visit the node.}
  \end{enumerate}%
}
\end{columns}
\end{frame}


\begin{frame}{Example of Binary Tree Traversal}
\begin{columns}

  \column{0.5\textwidth}
    \begin{itemize}
    \item
      \onslide<2->{%
        Input (one example):\\
        \texttt{\bfseries E C F A D H B G}%
      }
    \item
      In-order sequence:\\
      \onslide<4->{%
        \texttt{\bfseries A B C D E F G H}%
      }
    \item
      Pre-order sequence:\\
      \onslide<5->{%
        \texttt{\bfseries E C A B D F H G}\\
        (Use to copy and preserve a tree's structure.)%
      }
    \item
      Post-order sequence:\\
      \onslide<6->{%
      \texttt{\bfseries B A D C G H F E}\\
        (Use to delete a subtree.)%
      }
    \end{itemize}
  \column{0.5\textwidth}
    \centering
    \onslide<3->{
      \begin{forest}
      [E, for tree={fit=tight}
        [C
          [A,
            [\NullNode]
            [B
              [\NullNode]
              [\NullNode]
            ]
          ]
          [D
            [\NullNode]
            [\NullNode]
          ]
        ]
        [F
          [\NullNode]
          [H
            [G
              [\NullNode]
              [\NullNode]
            ]
            [\NullNode]
          ]
        ]
      ]
      \end{forest}%
    }
\end{columns}
\end{frame}

\begin{frame}{Depth- and Breadth-First Traversals}
  The in-order, pre-order, and post-order traversals we implemented are 
  \alt<1>{\uline{\invisible{depth}}}{\uline{depth}}-first traversals.

  \onslide<3->{We could also use a queue to implement a level-based (\emph{breadth-first}) traversal.}
\end{frame}

\section{Other Operations}
\begin{frame}{Typical Operations on Binary Tree}
\setbeamercovered{transparent}
\begin{itemize}[<+->]
\item
  Determine whether the binary tree is empty.
\item
  Insert an item in the binary tree.
\item
  Delete an item from the binary tree.
\item
  Search the binary tree for a particular item.
\item
  Find the height of the binary tree.
\item
  Find the number of nodes in the binary tree.
\item
  Find the number of leaves in the binary tree.
\item
  Traverse the binary tree.
\item
  Copy the binary tree.
\end{itemize}
\end{frame}

\section*{Lab 14}
\begin{frame}{Lab 14: Binary Trees (Part 1)}
\centering\Large
Let's take a look at the lab based on today's lecture.
\end{frame}
\end{document}
